\section{네이티브 앱 (macOS, iOS, Android, Desktop) 제작 방식}

ARES는 현재 웹 단독(HTML5 + Web Bluetooth + WebGL) 아키텍처로, 별도의
네이티브 앱 프레임워크(Electron, Capacitor 등)를 사용하지 않는다. 본 장은
기존 웹 자산을 최대한 재사용하면서 각 플랫폼용 네이티브 앱으로 패키징하는
방안을 제시한다.

\subsection{핵심 제약: 플랫폼별 BLE 지원 차이}

네이티브 패키징의 성패는 \textbf{Web Bluetooth 지원 여부}에 달려 있다.
현재 통신 계층(\code{bluetooth.js})은 \code{navigator.bluetooth}에 직접
의존하는데, 이 API의 지원은 플랫폼마다 크게 다르다.

\begin{longtable}{L{3.4cm} L{4.0cm} L{6.6cm}}
\toprule
\rowcolor{tablehead}\textbf{실행 환경} & \textbf{Web Bluetooth} & \textbf{함의} \\
\midrule
\endhead
데스크톱 Chrome/Edge & 지원 & 그대로 동작(현행 방식) \\
Electron(데스크톱) & 지원(핸들러 설정) & 웹 코드 거의 그대로 재사용 \\
Android Chrome & 지원 & 브라우저로는 동작 \\
Android WebView & 미지원 & 래퍼 앱은 네이티브 BLE 브리지 필요 \\
iOS Safari/WKWebView & 미지원 & 네이티브 BLE 브리지 필수 \\
\bottomrule
\end{longtable}

\begin{infobox}[title=설계 권고: 통신 계층 추상화]
네이티브화의 선행 작업으로, \code{bluetooth.js}를 \emph{전송(transport)
인터페이스}로 추상화할 것을 권한다. 즉 \code{connect()}, \code{sendData()},
\code{onReceive()} 등을 인터페이스로 정의하고, 구현체를 (1) Web Bluetooth,
(2) 네이티브 BLE 플러그인 브리지로 교체 가능하게 만든다. 명령 생성
(\code{commandexecutor.js})과 UI는 전송 구현에 무관하게 유지된다.
이 추상화 한 번으로 모든 플랫폼 분기를 흡수할 수 있다.
\end{infobox}

권고하는 전송 추상화의 연계 구조는 그림~\ref{fig:transport}와 같다. 상위
모듈(명령 생성\textbf{·}UI)은 \emph{전송 인터페이스}에만 의존하고, 실제 구현체는
플랫폼에 따라 교체된다. 현행 Web Bluetooth 구현은 데스크톱\textbf{·}Android
브라우저\textbf{·}Electron에서 그대로 쓰이고, WKWebView\textbf{·}Android WebView처럼
Web Bluetooth가 없는 환경에서는 네이티브 BLE 플러그인 구현으로 대체한다.

\begin{diagram}{전송 계층 추상화 \textbf{·} 플랫폼별 구현 연계도}{fig:transport}
\node[dom, text width=40mm] (app) {상위 모듈\\{\scriptsize commandexecutor · UI (불변)}};
\node[box, below=8mm of app, text width=52mm] (iface) {전송 인터페이스\\{\scriptsize connect() · sendData() · onReceive()}};
\node[hw, below left=10mm and 2mm of iface, text width=34mm] (web) {WebBluetoothTransport\\{\scriptsize navigator.bluetooth}};
\node[hw, below right=10mm and 2mm of iface, text width=34mm] (nat) {NativeBLETransport\\{\scriptsize Capacitor/CoreBluetooth 플러그인}};
\node[boxw, below=8mm of web, text width=34mm] (p1) {데스크톱·Electron\\Android 브라우저/PWA};
\node[boxw, below=8mm of nat, text width=34mm] (p2) {iOS·Android 래퍼\\{\scriptsize WebView(Web Bluetooth 없음)}};
\draw[dep] (app) -- node[lbl,right]{의존} (iface);
\draw[flow] (iface) -- node[lbl,left]{구현} (web);
\draw[flow] (iface) -- node[lbl,right]{구현} (nat);
\draw[flow] (web) -- (p1);
\draw[flow] (nat) -- (p2);
\end{diagram}

\subsection{데스크톱 (macOS / Windows / Linux)}

\subsubsection{방안 A --- Electron (권장)}

Electron\footnote{\textbf{Electron}: 웹 기술(HTML/CSS/JS)로 데스크톱 앱을 만드는 프레임워크.
Chromium 브라우저 엔진과 Node.js를 내장한다.}은 Chromium\footnote{\textbf{Chromium}: Google
Chrome\textbf{·}Microsoft Edge 등의 기반이 되는 오픈소스 브라우저 엔진.}을 내장하므로
Web Bluetooth를 지원하며, 기존
\code{Build/} 산출물을 거의 수정 없이 탑재할 수 있다.

\begin{itemize}
  \item 메인 프로세스에서 \code{select-bluetooth-device} 이벤트를 처리하여
        장치 선택 UI를 연결한다(Electron은 기본 브라우저식 장치 선택기를
        제공하지 않으므로 핸들러가 필요).
  \item 오프라인 빌드의 fetch 심\textbf{·}로컬 vendor 자산을 그대로 활용한다.
  \item \code{electron-builder}로 macOS \code{.dmg}/\code{.app},
        Windows \code{.exe}/\code{installer}, Linux \code{AppImage} 생성.
  \item macOS는 BLE 사용을 위해 \code{NSBluetoothAlwaysUsageDescription} 권한
        고지와 코드 서명\textbf{·}공증(notarization)이 필요하다.
\end{itemize}

\subsubsection{방안 B --- Tauri (경량 대안)}

Tauri\footnote{\textbf{Tauri}: 시스템에 내장된 WebView와 Rust 백엔드를 사용해 가벼운 데스크톱
앱을 만드는 프레임워크.}는 시스템 WebView + Rust 백엔드로 번들 크기가 작다. 다만 macOS의
WKWebView\footnote{\textbf{WKWebView}: iOS\textbf{·}macOS 앱에 웹 콘텐츠를 표시하는 애플의 웹뷰
컴포넌트. Web Bluetooth를 지원하지 않는다.}는 Web Bluetooth를 지원하지 않으므로, Rust 측 BLE 라이브러리
(예: \code{btleplug})로 네이티브 BLE를 구현하고 JS와 브리지해야 한다.
위의 전송 추상화가 전제된다.

\subsection{Android}

\begin{itemize}
  \item \textbf{Capacitor\footnote{\textbf{Capacitor}: 웹 앱을 iOS\textbf{·}Android 네이티브 앱으로
        감싸고, 카메라\textbf{·}블루투스 같은 네이티브 기능을 플러그인으로 연결해 주는 런타임.} 래퍼}:
        \code{Build/} 웹 자산을 Capacitor 프로젝트의
        웹 디렉터리로 사용한다. 단, Android System WebView는 Web Bluetooth를
        지원하지 않으므로, BLE는 네이티브 플러그인(예:
        \code{@capacitor-community/bluetooth-le})으로 처리하고, 전송 구현체를
        이 플러그인 호출로 교체한다.
  \item \textbf{권한}: \code{BLUETOOTH\_SCAN}, \code{BLUETOOTH\_CONNECT}
        (Android 12+), 구형 기기 스캔을 위한 위치 권한을 매니페스트에 선언한다.
  \item \textbf{대안(PWA\footnote{\textbf{PWA(Progressive Web App)}: 브라우저로 접속하지만 홈 화면
        설치\textbf{·}오프라인 사용\textbf{·}푸시 알림 등 네이티브 앱에 가까운 경험을 제공하는 웹 앱.})}:
        Android Chrome은 Web Bluetooth를 지원하므로, 설치형
        PWA로 배포하면 네이티브 빌드 없이 홈 화면 설치\textbf{·}오프라인 캐시가
        가능하다.
\end{itemize}

\subsection{iOS}

iOS는 Safari\textbf{·}WKWebView 모두 Web Bluetooth를 지원하지 않는다. 따라서
순수 웹 방식으로는 BLE 제어가 불가능하며, 다음 중 하나가 필요하다.

\begin{itemize}
  \item \textbf{Capacitor + 네이티브 BLE 플러그인}: 웹 UI는 재사용하되, BLE는
        CoreBluetooth 기반 플러그인(\code{@capacitor-community/bluetooth-le})으로
        구현하고 전송 추상화로 연결한다. App Store 배포 시
        \code{NSBluetoothAlwaysUsageDescription} 고지가 필수이다.
  \item \textbf{Web Bluetooth 지원 브라우저(WebBLE/Bluefy)}: 별도 앱 없이
        해당 서드파티 브라우저로 웹앱을 여는 임시 방안. 배포\textbf{·}통제 측면에서
        교육 현장 권장도는 낮다.
\end{itemize}

\subsection{플랫폼별 권장 조합 요약}

\begin{longtable}{L{2.8cm} L{3.6cm} L{7.6cm}}
\toprule
\rowcolor{tablehead}\textbf{플랫폼} & \textbf{권장 방식} & \textbf{BLE 처리} \\
\midrule
\endhead
macOS/Win/Linux & Electron & Web Bluetooth(핸들러 설정) --- 코드 재사용 최대 \\
Android & PWA 또는 Capacitor & PWA는 Web Bluetooth, 래퍼는 네이티브 BLE 플러그인 \\
iOS & Capacitor & 네이티브 BLE 플러그인(CoreBluetooth) 필수 \\
공통 선행작업 & 전송 추상화 & \code{bluetooth.js}를 인터페이스화 \\
\bottomrule
\end{longtable}

\subsection{단계적 도입 로드맵}

\begin{enumerate}
  \item \code{bluetooth.js}를 전송 인터페이스로 리팩터링(Web Bluetooth 구현을
        기본 구현체로 유지).
  \item Electron 데스크톱 앱부터 출시(코드 재사용도가 가장 높음).
  \item Android는 PWA로 우선 제공, 필요 시 Capacitor + BLE 플러그인으로 승격.
  \item iOS는 Capacitor + 네이티브 BLE 플러그인으로 구현(가장 큰 추가 작업).
  \item 코드 서명\textbf{·}공증\textbf{·}스토어 심사\textbf{·}권한 고지 등 배포 요건을
        플랫폼별로 정리.
\end{enumerate}

\begin{teachbox}{한 번에 여러 플랫폼으로}
\teachhead{설계 의도}
\begin{itemize}
  \item 플랫폼마다 다른 BLE 지원 차이를 \emph{전송 추상화} 한 곳에서 흡수해, 상위 코드를 건드리지 않는다.
\end{itemize}
\teachhead{분석할 때 핵심 질문}
\begin{itemize}
  \item 왜 iOS는 웹만으로 BLE 제어가 안 되는가?
  \item 전송 추상화를 먼저 해 두면 네이티브 확장이 왜 쉬워지는가?
\end{itemize}
\teachhead{점검 요소}
\begin{itemize}
  \item \code{bluetooth.js}를 \code{connect}/\code{sendData}/\code{onReceive} 인터페이스로 분리하는 설계를 종이에 그려 보라.
\end{itemize}
\end{teachbox}

